\chapter*{前言}
\addcontentsline{toc}{chapter}{前言}
\markboth{前言}{前言}

本书基于「模空间的几何与量子化」（The Geometry and Quantization of Moduli Spaces）
学期课程期间开设的四门高级课程，该学期课程于 2012 年 3 月至 6 月在巴塞罗那
Bellaterra 的数学研究中心（Centre de Recerca Matemàtica，CRM）举行。模空间除了
在数学的许多领域中扮演重要角色之外，近几十年来还成为理解高能物理中若干现象的关键，
并因此促成了数学与物理之间的互动——这种互动对两门学科而言都可谓收获惊人。我们希望
读者能在本卷中欣赏并享受这种互动的美。

这些课程所处理的几个专题，在模空间这一广阔研究领域内正经历非同寻常的增长，它们也
反映了该研究学期中反复出现的一条线索，即黎曼面上局部系统（local system）的模空间，
无论是从经典的角度还是从量子的角度。

第 1 章由 V.~V.~Fock（法国斯特拉斯堡，IRMA）与 A.~Marshakov（俄罗斯莫斯科，ITEP）
撰写，题为《圈群、cluster、二聚体与可积系统》（Loop Groups, Clusters, Dimers and
Integrable Systems），属于 cluster 可积系统的范畴；更确切地说，属于 A.~B.~Goncharov
与 R.~Kenyon 借助二维环面上二部图的统计物理中二聚体模型（dimer model）理论所构造的
那一类可积系统。这篇研究论文描述了这些可积系统与在适当的仿射 Poisson--Lie 群上构造的
标准可积系统之间的一个同构，从而既给出了前者的新构造，也给出了后者的新视角。本章包含
非常有趣的例子，并指出了若干令人兴奋的推广方向。

第 2 章由 F.~Schaffhauser（哥伦比亚波哥大，Universidad de los Andes）撰写，题为
《克莱因曲面及其基本群讲义》（Lectures on Klein Surfaces and Their Fundamental
Group），是关于克莱因曲面的基本群的一篇综述性讲义。它包括克莱因曲面与实代数曲线的
入门、对克莱因曲面基本群的细致处理，以及对克莱因曲面基本群的线性表示与酉表示、
酉表示簇的讨论。许多定义与结果都以简单而有描述力的例子作动机或说明，并且每一节都配有
有趣的习题，读者可用以巩固所学。

接下来是第 3 章，由 C.~Teleman（英国牛津，牛津大学）撰写，题为《拓扑场论五讲》
（Five Lectures on Topological Field Theory）。它从 M.~F.~Atiyah 整理成体系的
拓扑量子场论（TQFT）经典定义出发，经由不同的例子，引出这些理论如今在多数情形下已知
具有的额外结构——正如 J.~Lurie 最近证明的所谓配边猜想（cobordism hypothesis）所
预言的那样。作者并未给出扩展 TQFT 的完全详细的定义（读者可在 Lurie 的论文中找到），
而是借助精心挑选的例子，阐释这一观念实质的各个侧面，以及证明其相容性与配边猜想所面临
的困难。在我们看来，这使本文成为 Lurie 论文的一份极佳 companion。
\begin{annote}
「companion」在原文中即英文词，指与之配套、可对读的著作。此处照录，未译作「伴侣」。
\end{annote}

第四也是最后一章，由 R.~A.~Wentworth（美国 College Park，马里兰大学）撰写，题为
《黎曼面上的 Higgs 丛与局部系统》（Higgs Bundles and Local Systems on Riemann
Surfaces）。Higgs 丛由 N.~Hitchin 在三十年前引入，现已在不同领域中证明具有中心地位，
例如低维拓扑、代数几何、数学物理，乃至数论（比如几年前 B.~C.~Ngô 对 Langlands 纲领中
基本引理的证明）。本章对这一概念作了教学式的入门介绍，但限于黎曼面的情形（这正是
Hitchin 论文中的原始语境）。它给出了 Higgs 丛模空间与局部系统模空间之间对应的一个
近乎自足的证明，然后将其应用于研究 oper 模空间——这一对象在几何 Langlands 纲领近来的
发展中起着关键作用。

我们要感谢本书各位作者的重大贡献，既为准备了一套非常有趣的讲义，也为撰写了本册笔记。
我们也感谢各位审稿人，他们毫不妥协地帮助了我们。

「模空间的几何与量子化」学期课程得以举办，有赖于不同机构的慷慨资助。为此我们要感谢
数学研究中心（CRM）、通过研究网络 ITGP（Interactions between Topology, Geometry
and Physics，拓扑、几何与物理之间的相互作用）给予支持的欧洲研究基金会、IMUB
（Institut de Matemàtiques de la Universitat de Barcelona，巴塞罗那大学数学研究所）
以及 SCM（Societat Catalana de Matemàtiques，加泰罗尼亚数学学会）。最后，我们感谢
CRM 的工作人员，他们在课程筹备与运行期间的效率与帮助令人感念。

\vspace{1.2em}
\begin{flushright}
马德里，Luis Álvarez-Cónsul（ICMAT）\\
奥胡斯，Jørgen Ellegaard Andersen（QGM）\\
巴塞罗那，Ignasi Mundet i Riera（UB）
\end{flushright}
